(1) 磁准静态 Maxwell 方程 $\nabla\times\boldsymbol H=\boldsymbol J=\sigma\boldsymbol E$，$\nabla\times\boldsymbol E=-\partial_t\boldsymbol B$，$\boldsymbol B=\mu\boldsymbol H$。在均匀导体中取旋度并用 $\nabla\cdot\boldsymbol B=0$，得到 $\nabla^2\boldsymbol B=\mu\sigma\partial_t\boldsymbol B$。设场仅沿深度 $x$ 变化，便有 $B_{xx}=\mu\sigma B_t$。
(2) 用复表示 $H_y(x,t)=\Re[H_0e^{-kx}e^{-i\omega t}]$，代入给出 $k^2=-i\omega\mu\sigma$；取 $\Re k>0$ 分支可写 $k=(1-i)/\delta$，$\delta=\sqrt{2/(\mu\sigma\omega)}$。因此实际场为 $H_y=H_0e^{-x/\delta}\cos(\omega t-x/\delta)$（若 $H_0$ 为实）。场幅沿导体深度衰减，传播相位也随深度延迟。
(3) 由 $\boldsymbol J=\nabla\times\boldsymbol H$，切向电流密度复幅大小为 $|J|=|kH_0|e^{-x/\delta}=\sqrt2|H_0|e^{-x/\delta}/\delta$。一个周期的平均焦耳热体功率为 $\langle J^2/\sigma\rangle=|H_0|^2e^{-2x/\delta}/(\sigma\delta^2)$。积分 $x\ge0$ 得 $P_{\rm heat}/A=|H_0|^2/(2\sigma\delta)$，说明在高频时耗散集中于薄表层。
(4) 定义表面电阻 $R_s=1/(\sigma\delta)=\sqrt{\mu\omega/(2\sigma)}$，则单位面积耗散率为 $R_s|H_0|^2/2$。$\omega$ 加倍时 $\delta$ 变为原来的 $1/\sqrt2$、$R_s$ 增大 $\sqrt2$。$\sigma\to\infty$ 时理想模型 $\delta\to0$、$R_s\to0$，但若出现反常集肤效应、非局域输运或忽略位移电流失效，就需改用更完整的材料响应。
